\sloppy
\emergencystretch=5em
\hbadness=10000
\vbadness=10000
\tolerance=2000

\providecolor{codebg}{RGB}{248,248,248}
\providecolor{codeframe}{RGB}{200,200,200}
\providecolor{keyword}{RGB}{0,0,180}
\providecolor{string}{RGB}{163,21,21}
\providecolor{comment}{RGB}{0,128,0}

\lstset{
	backgroundcolor=\color{codebg},
	frame=single,
	rulecolor=\color{codeframe},
	language=Python,
	basicstyle=\ttfamily\scriptsize,
	keywordstyle=\color{keyword}\bfseries,
	stringstyle=\color{string},
	commentstyle=\color{comment}\itshape,
	showstringspaces=false,
	breaklines=true,
	breakatwhitespace=true,
	numbers=left,
	numberstyle=\tiny\color{gray},
	numbersep=8pt,
	tabsize=4,
	literate=
	{ف}{{{\bfseries\itshape ف}}}1
	{ی}{{{\bfseries\itshape ی}}}1
}

\section{نمای کلی و هدف ماژول}

ماژول \lr{\texttt{validate\_physics.py}} یک اسکریپت اعتبارسنجی فیزیکی است که نقش \emph{بین‌اللایه‌ای حیاتی} در فریم‌ورک شبیه‌سازی \lr{QKD} ایفا می‌کند. این ماژول به‌عنوان یک \lr{golden reference} عمل می‌کند که خروجی موتور شبیه‌سازی کامل \lr{\texttt{main\_optimized.py}} را در برابر معادلات تحلیلی بسته و مستقل از پیاده‌سازی، مقایسه می‌نماید و با استفاده از آزمون آماری \lr{Z-test} دودویی (\lr{binomial Z-test}) در سطح معنی‌داری $\alpha = 10^{-4}$، انطباق آماری نتایج را بررسی می‌سازد. وجود چنین اسکریپت اعتبارسنجی جداگانه‌ای در یک فریم‌ورک علمی، نشانه بلوغ مهندسی نرم‌افزار است؛ زیرا هرگونه تغییر در پیاده‌سازی موتور فیزیکی، چه عمدی (مانند اصلاح یک باگ) و چه ناخواسته (مانند \lr{regression} پس از \lr{refactor})، بلافاصله توسط این آزمون‌ها شناسایی می‌شود و پیش از انتشار نتایج، به اطلاع توسعه‌دهنده می‌رسد.

هدف اصلی این ماژول، پاسخ به پرسش بنیادین زیر است:
\emph{آیا خروجی شبیه‌سازی مونت‌کارلو که از خط لولهٔ کامل
	\lr{\texttt{Source -> Channel -> Protocol -> Detector -> Sifting -> Proof}}
	به دست می‌آید، از نظر آماری با پیش‌بینی تئوریکِ معادلات تحلیلیِ بسته هم‌خوان است؟}
این پرسش به دو دلیل اهمیت دارد. نخست، در شبیه‌سازی‌های فیزیکی که در آن‌ها فرایندهای تصادفی متعددی (نظیر توزیع \lr{Poisson} فوتون‌ها، شمارش تاریکِ آشکارساز (\lr{dark count})، غربال‌سازیِ پایه‌ها (\lr{sifting}) و غیره) با یکدیگر ترکیب می‌شوند، خطاهای ظریف در پیاده‌سازی ممکن است به نتایجی منجر شود که از نظر فیزیکی قابل قبول به نظر می‌رسند، اما از نظر آماری \emph{biased} هستند. دوم، معادلات تحلیلیِ بسته که در ادبیات \lr{QKD} به‌طور گسترده پذیرفته شده‌اند (مانند معادلات \lr{Gain} و \lr{QBER} در مدل \lr{Gottesman--Lo--Lütkenhaus--Preskill}) به‌عنوان مرجع علمی عمل می‌کنند و انطباق با آن‌ها، اعتبار فیزیکی پیاده‌سازی را تضمین می‌نماید.


ساختار کلی ماژول شامل چهار تابع محاسباتی تحلیلی، یک تابع آزمون آماری و یک تابع اجرای شبیه‌سازی ایزوله است. تابع \lr{\texttt{channel\_transmittance}} ضریب انتقال کانال فیبر نوری را بر اساس قانون \lr{Beer-Lambert} محاسبه می‌کند. تابع \lr{\texttt{Q\_mu\_poisson}} \lr{Gain} (نرخ کلیک به ازای هر پالس) را برای منبع \lr{Weak Coherent Pulse} با توزیع \lr{Poisson} محاسبه می‌کند. تابع \lr{\texttt{Q\_mu\_thermal}} همین کمیت را برای منبع \lr{Thermal} (تک‌حالت یا چندحالت) محاسبه می‌نماید. تابع \lr{\texttt{E\_mu\_poisson}} \lr{QBER} متوسط را برای منبع \lr{Poisson} به دست می‌آورد. تابع \lr{\texttt{dead\_time\_correction}} تصحیح \lr{dead time} آشکارساز را بر اساس مدل \lr{Non-Paralyzable} اعمال می‌کند. تابع \lr{\texttt{validate\_statistic}} آزمون \lr{Z-test} دودویی را روی یک مشاهده انجام می‌دهد و نتیجه \lr{PASS}/\lr{FAIL} را برمی‌گرداند. در نهایت، تابع \lr{\texttt{run\_isolated\_simulation}} یک نشست شبیه‌سازی مستقل روی موتور \lr{\texttt{main\_optimized.py}} اجرا می‌کند و خروجی را به فرمت مناسب برای آزمون آماری بازمی‌گرداند.

سه گروه آزمون در این ماژول تعریف شده است که هر یک به جنبه‌ای متفاوت از فیزیک شبیه‌سازی می‌پردازد. \lr{TEST 1} منبع \lr{Weak Coherent Pulse} با توزیع \lr{Poisson} را در سه فاصله $0$، $50$ و $100$ کیلومتر اعتبارسنجی می‌کند؛ این آزمون هم \lr{Gain} و هم \lr{QBER} را در رژیم مجانبی (یعنی بدون \lr{decoy} و بدون \lr{dead time}) بررسی می‌سازد. \lr{TEST 2} منبع \lr{Thermal} تک‌حالت را در فاصله صفر آزمون می‌کند تا توزیع آماری متفاوت منبع را راستی‌آزمایی نماید. \lr{TEST 3} فیزیک \lr{dead time} آشکارساز \lr{SPAD} را در حضور دو آشکارساز مستقل بررسی می‌کند و این موضوع مهم را که هر \lr{SPAD} به‌طور مستقل زمان مرده خود را تجربه می‌کند، به آزمون می‌گذارد.

یکی از ویژگی‌های برجسته این ماژول، انتخاب سطح معنی‌داری $\alpha = 10^{-4}$ است که به \lr{Z-crit} $\approx 3.8906$ متناظر می‌شود. این سطح به‌مراتب سخت‌گیرانه‌تر از سطح متداول $\alpha = 0.05$ (\lr{Z-crit} $\approx 1.96$) یا حتی $\alpha = 0.001$ (\lr{Z-crit} $\approx 3.29$) است. دلیل این انتخاب سخت‌گیرانه آن است که در شبیه‌سازی \lr{QKD} با $N = 100000$ پالس، حتی انحرافات کوچک از تئوری می‌توانند به‌سرعت شناسایی شوند و توسعه‌دهنده را از پذیرش نتایج \emph{تقریباً درست} که در عمل \emph{biased} هستند، بازدارند. با این حال، این سطح به‌قدری سخت‌گیرانه نیست که نویز آماری ذاتی شبیه‌سازی مونت‌کارلو باعث \lr{false positive} مکرر شود؛ تعادل میان حساسیت و \lr{false positive rate} با دقت انتخاب شده است.

در سطح عملیاتی، این ماژول به‌عنوان یک \lr{CI/CD gate} عمل می‌کند؛ یعنی پیش از \lr{merge} هر تغییر در شاخه اصلی، این اسکریپت اجرا می‌شود و در صورت شکست هر یک از آزمون‌ها، \lr{pipeline} متوقف می‌گردد. این رویکرد \lr{shift-left testing} تضمین می‌کند که باگ‌های فیزیکی در همان مراحل اولیه توسعه شناسایی شوند، نه پس از آنکه نتایج شبیه‌سازی در مقالات یا گزارش‌های فنی منتشر شده باشند. در ادبیات مهندسی نرم‌افزار علمی، این الگو به‌عنوان \lr{regression testing for computational physics} شناخته می‌شود و در پروژه‌های بالغی مانند \lr{Quantum ESPRESSO}، \lr{LAMMPS} و \lr{Qiskit} به‌صورت استاندارد به کار می‌رود.

\section{مبانی تئوری و فیزیکی}

\subsection{کانال فیبر نوری و قانون Beer-Lambert}

انتقال فوتون در فیبر نوری تک‌حالت با قانون \lr{Beer-Lambert} مدل می‌شود که رابطه بنیادین بین تلفات فیبر، فاصله و ضریب انتقال را بیان می‌دارد. برای یک فیبر به طول $L$ کیلومتر با ضریب تلفت خطی $\alpha_{\text{dB/km}}$ (بر حسب دسی‌بل بر کیلومتر)، تلفات کلی کانال بر حسب دسی‌بل برابر با $\alpha_{\text{dB/km}} \cdot L$ است. چون مقیاس دسی‌بل لگاریتمی است، تبدیل آن به ضریب انتقال خطی (\lr{power transmittance}) از رابطه زیر پیروی می‌کند:

\begin{equation}
	\eta_{\text{channel}}(L, \alpha_{\text{dB/km}}) = 10^{-\frac{\alpha_{\text{dB/km}} \cdot L}{10}}
\end{equation}

در این رابطه، $L$ بر حسب کیلومتر و $\alpha_{\text{dB/km}}$ بر حسب \lr{dB/km} است. در طول‌موج مخابراتی استاندارد $1550$ نانومتر، فیبر تک‌حالت استاندارد ضریب تلفت $\alpha_{\text{dB/km}} \approx 0.2$ دارد؛ بنابراین در فاصله $50$ کیلومتر، ضریب انتقال کانال برابر با $10^{-1} = 0.1$ و در فاصله $100$ کیلومتر، برابر با $10^{-2} = 0.01$ می‌شود. این کاهش نمایی ضریب انتقال با فاصله، عامل اصلی محدودیت برد \lr{QKD} در فیبر نوری است و دلیل آن است که در فواصل بزرگ‌تر از $200$ کیلومتر، حتی با استفاده از \lr{decoy-state}، نرخ کلید امن به صفر میل می‌کند.

ضریب انتقال کلی سیستم شامل بازده آشکارساز $\eta_{\text{det}}$ نیز می‌شود که در این ماژول به‌صورت یک کمیت مستقل در نظر گرفته شده است:

\begin{equation}
	\eta_{\text{total}}(L) = \eta_{\text{det}} \cdot \eta_{\text{channel}}(L, \alpha_{\text{dB/km}})
\end{equation}

در پیاده‌سازی پیش‌فرض این ماژول، $\eta_{\text{det}} = 0.15$ است که متناظر با آشکارساز \lr{InGaAs SPAD} در دمای $-50$ درجه سانتی‌گراد با بازده کوانتومی متداول است. در عمل، این مقدار در بازه $0.1$ تا $0.25$ متغیر است و به دمای آشکارساز، ولتاژ بایاس و طول‌موج فوتون بستگی دارد.

\subsection{منبع Weak Coherent Pulse و توزیع Poisson}

در اکثر پیاده‌سازی‌های عملی \lr{QKD}، به‌جای منبع تک‌فوتونی آرمانی (که هنوز در دسترس تجاری نیست)، از منبع \lr{Weak Coherent Pulse} (\lr{WCP}) استفاده می‌شود. در این منبع، لیزر با شدت متوسط $\mu$ فوتون در هر پالس تحریک می‌شود و تعداد فوتون‌های هر پالس از توزیع \lr{Poisson} پیروی می‌کند:

\begin{equation}
	P(n | \mu) = \frac{\mu^n e^{-\mu}}{n!}, \quad n = 0, 1, 2, \ldots
\end{equation}

با این توزیع، احتمال اینکه یک پالس حاوی $n$ فوتون باشد به‌صورت بالا داده می‌شود. پارامتر $\mu$ میانگین تعداد فوتون در هر پالس است و در \lr{decoy-state BB84} معمولاً در بازه $0.1$ تا $0.5$ انتخاب می‌شود. مقدار $\mu = 0.5$ که در این ماژول به‌کار رفته، یک انتخاب متعادل میان نرخ کلید (که با $\mu$ رشد می‌کند) و امنیت (که با $\mu$ کاهش می‌یابد، چون پالس‌های چندفوتونی به حمله \lr{Photon-Number Splitting} حساس هستند) است.

\lr{Gain} کل سیستم (\lr{detection yield}) که با $Q_\mu$ نشان داده می‌شود، احتمال آن است که یک پالس منجر به ثبت کلیک در آشکارساز شود. برای منبع \lr{WCP} با توزیع \lr{Poisson}، در حضور دارک کانت با نرخ $Y_0$ و بازده کلی $\eta$، این کمیت از رابطه بسته زیر به دست می‌آید:

\begin{equation}
	Q_\mu^{\text{Poisson}}(\mu, \eta, Y_0) = 1 - (1 - Y_0) \cdot e^{-\mu \cdot \eta}
\end{equation}

این رابطه از جمع احتمالات تمام حالات ممکن به دست می‌آید. به‌طور شهودی، عبارت $e^{-\mu \cdot \eta}$ احتمال آن است که هیچ فوتونی شناسایی نشود (یعنی یا هیچ فوتونی ارسال نشده، یا فوتون‌های ارسالی از دست رفته‌اند) و عامل $(1 - Y_0)$ احتمال آن است که دارک کانت نیز رخ ندهد. بنابراین، $1 - (1-Y_0) e^{-\mu \eta}$ احتمال متمم است: یعنی احتمال آنکه حداقل یک کلیک (چه از فوتون واقعی و چه از دارک کانت) ثبت شود. این فرمول در ادبیات \lr{QKD} به‌ویژه در کار کلاسیک \lr{Lo-Ma-Chen (2005)} به‌تفصیل استخراج شده است.

\subsection{منبع Thermal و توزیع Bose-Einstein}

در برخی سناریوهای \lr{QKD}، به‌جای لیزر \lr{WCP} از منبع \lr{Thermal} استفاده می‌شود. منبع حرارتی فوتون‌ها با توزیع \lr{Bose-Einstein} تولید می‌کند که برای یک منبع تک‌حالتی به‌صورت زیر تعریف می‌شود:

\begin{equation}
	P(n | \mu) = \frac{\mu^n}{(1 + \mu)^{n+1}}, \quad n = 0, 1, 2, \ldots
\end{equation}

برای یک منبع چندحالتی با $M$ حالت مستقل، توزیع مرجع به توزیع \lr{Negative Binomial} میل می‌کند. \lr{Gain} کل سیستم برای منبع حرارتی با $M$ حالت، از رابطه زیر به دست می‌آید:

\begin{equation}
	Q_\mu^{\text{Thermal}}(\mu, \eta, Y_0, M) = 1 - \frac{1 - Y_0}{\left(1 + \frac{\mu \cdot \eta}{M}\right)^M}
\end{equation}

این رابطه تعمیم طبیعی فرمول \lr{Poisson} است. در حد $M \to \infty$، توزیع \lr{Thermal} به \lr{Poisson} میل می‌کند و فرمول بالا به فرمول \lr{Poisson} تقلیل می‌یابد. برای $M = 1$ (تک‌حالت)، این رابطه به $Q_\mu^{\text{Thermal}} = 1 - (1-Y_0)/(1+\mu \eta)$ ساده می‌شود که فرمول کلاسیک منبع حرارتی تک‌حالت است. تفاوت فیزیکی میان \lr{Poisson} و \lr{Thermal} در آن است که در \lr{Thermal}، فوتون‌ها تمایل به \emph{bunching} دارند (یعنی احتمال پالس‌های چندفوتونی بیشتر است) که این امر امنیت را در برابر حمله \lr{PNS} کاهش می‌دهد. با این حال، مطالعه \lr{Thermal} از منظر آکادمیک اهمیت دارد چرا که برخی منابع نور تک‌فوتونی (مانند \lr{quantum dot}) در رژیم خاصی رفتار حرارتی نشان می‌دهند.

\subsection{مدل QBER برای منبع Poisson}

نرخ خطای کوانتومی بیت (\lr{Quantum Bit Error Rate}) که با $E_\mu$ نشان داده می‌شود، کسری از بیت‌های \lr{sifted} است که بین دو طرف ناهماهنگ هستند. در مدل تحلیلی، این کمیت از نسبت خطاهای کلیک به کل کلیک‌ها به دست می‌آید:

\begin{equation}
	E_\mu^{\text{Poisson}}(\mu, \eta, Y_0, e_d) = \frac{0.5 \cdot Y_0 + e_d \cdot (Q_\mu - Y_0)}{Q_\mu}
\end{equation}

در این رابطه، $e_d$ خطای نوری ذاتی (\lr{misalignment error}) است که از نقص‌های تداخل‌سنج ناشی می‌شود و $Y_0$ نرخ دارک کانت است. صورت این کسر نشان می‌دهد که خطا از دو منبع می‌آید: دارک کانت که کاملاً تصادفی است و نرخ خطای $0.5$ دارد (یعنی نیمی از دارک کانت‌ها با بیت ارسالی هم‌خوان و نیمی ناهم‌خوان هستند)، و فوتون‌های واقعی که با احتمال $e_d$ دچار خطا می‌شوند. مخرج $Q_\mu$ کل کلیک‌ها است که مجموع دارک و سیگنال است. در نتیجه، $E_\mu$ کسر وزنی خطاها است.

این فرمول نشان می‌دهد که در حضور دارک کانت، حتی اگر $e_d = 0$ باشد (یعنی هیچ خطای نوری ذاتی وجود نداشته باشد)، \lr{QBER} همچنان بزرگ‌تر از صفر خواهد بود؛ این پدیده به‌ویژه در فواصل طولانی که سیگنال ضعیف می‌شود، اهمیت پیدا می‌کند و دلیل اصلی آن است که در فواصل بزرگ، \lr{QBER} به $0.5$ میل می‌کند (زیرا بیشتر کلیک‌ها دارک هستند). در حد $\eta \to 0$ (یعنی کانال با تلفات بسیار بالا)، $Q_\mu \to Y_0$ و در نتیجه $E_\mu \to 0.5$ که نشان‌دهنده از دست رفتن کامل اطلاعات است.

در صورت کسر، عبارت $0.5 \cdot Y_0$ نشان‌دهنده آن است که دارک کانت‌ها به‌صورت کاملاً تصادفی رخ می‌دهند و بنابراین نرخ خطای آن‌ها برابر $0.5$ است. عبارت $e_d \cdot (Q_\mu - Y_0)$ نشان‌دهنده خطای ناشی از فوتون‌های واقعی است؛ $(Q_\mu - Y_0)$ بخشی از کلیک‌ها را نشان می‌دهد که از فوتون‌های واقعی می‌آیند (نه از دارک کانت) و این بخش با احتمال $e_d$ دچار خطا می‌شود. بنابراین، صورت کل کسر برابر است با کل خطاهای ثبت‌شده در واحد پالس.

\subsection{مدل Dead Time آشکارساز SPAD}

آشکارسازهای \lr{Single-Photon Avalanche Diode} (\lr{SPAD}) پس از هر کلیک، به مدتی برای بازیابی نیاز دارند که به آن \lr{dead time} می‌گویند. در این بازه زمانی $\tau_d$، آشکارساز قادر به ثبت کلیک جدیدی نیست. دو مدل متداول برای این پدیده وجود دارد: \lr{Non-Paralyzable} و \lr{Paralyzable}. در این ماژول، مدل \lr{Non-Paralyzable} به‌کار رفته است که در آن، کلیک‌های رخ‌داده در زمان مرده آشکارساز را برای مدت بیشتر \emph{بازنشانی نمی‌کنند} بلکه صرفاً نادیده گرفته می‌شوند.

برای یک آشکارساز \lr{Non-Paralyzable} با احتمال کلیک ایده‌آل $Q_{\text{ideal}}$ به ازای هر پالس و بازه زمانی میان پالس‌ها $T_{\text{pulse}}$، احتمال کلیک مشاهده‌شده $Q_{\text{obs}}$ از رابطه زیر به دست می‌آید:

\begin{equation}
	Q_{\text{obs}} = \frac{Q_{\text{ideal}}}{1 + Q_{\text{ideal}} \cdot \frac{\tau_d}{T_{\text{pulse}}}}
\end{equation}

این رابطه از حل معادله تابع ضمنی $Q_{\text{obs}} = Q_{\text{ideal}} \cdot (1 - Q_{\text{obs}} \cdot \tau_d / T_{\text{pulse}})$ به دست می‌آید که نشان می‌دهد کسر زمانی که آشکارساز در دسترس است، برابر با $1 - Q_{\text{obs}} \cdot \tau_d / T_{\text{pulse}}$ است. در حد $\tau_d \to 0$، $Q_{\text{obs}} \to Q_{\text{ideal}}$ (بدون کاهش) و در حد $\tau_d \to \infty$ یا $Q_{\text{ideal}} \to 1$، $Q_{\text{obs}} \to T_{\text{pulse}} / \tau_d$ (اشباع).

یکی از نکات ظریف در پیاده‌سازی این ماژول، آن است که در سیستم \lr{BB84} دو آشکارساز مستقل \lr{D0} و \lr{D1} وجود دارند که هر کدام زمان مرده \emph{مستقل} خود را تجربه می‌کنند. بنابراین، احتمال کلیک \lr{Gain} ایده‌آل ابتدا باید به دو نیمه تقسیم شود (هر آشکارساز نیمی از کلیک‌ها را ثبت می‌کند)، سپس تصحیح \lr{dead time} به‌صورت مستقل روی هر آشکارساز اعمال شود:

\begin{equation}
	Q_{\text{per-spad}} = \frac{Q_{\text{ideal}}^{\text{raw}}}{2}, \quad Q_{\text{obs}}^{\text{per-spad}} = \frac{Q_{\text{per-spad}}}{1 + Q_{\text{per-spad}} \cdot \tau_d / T_{\text{pulse}}}
\end{equation}

\begin{equation}
	Q_{\text{obs}}^{\text{total}} = 2 \cdot Q_{\text{obs}}^{\text{per-spad}}
\end{equation}

در نهایت، چون \lr{sifting} در \lr{BB84} نیمی از پالس‌ها را حذف می‌کند (پالس‌هایی که پایه گیرنده با فرستنده هم‌خوان نداشته)، \lr{Gain} نهایی که در خروجی شبیه‌سازی مشاهده می‌شود، برابر با $0.5 \cdot Q_{\text{obs}}^{\text{total}}$ است:

\begin{equation}
	Q_{\text{final}} = Q_{\text{obs}}^{\text{per-spad}} \cdot 2 \cdot 0.5 = Q_{\text{obs}}^{\text{per-spad}}
\end{equation}

این رابطه ظریف نشان می‌دهد که در نهایت، \lr{Gain} مشاهده‌شده با \lr{Gain} هر آشکارساز برابر است؛ چیزی که در نگاه اول غیرمنتظره به نظر می‌رسد اما با بررسی دقیق عوامل $2$ و $0.5$ که یکدیگر را خنثی می‌کنند، آشکار می‌گردد.

\subsection{آزمون آماری Z-test دودویی}

برای اعتبارسنجی آماری نتایج شبیه‌سازی، این ماژول از \lr{Z-test} دودویی استفاده می‌کند که یکی از ساده‌ترین و در عین حال قدرتمندترین آزمون‌های آماری برای متغیرهای دودویی است. فرض کنید $N$ آزمایش مستقل انجام شده و $k$ مورد موفقیت مشاهده شده است؛ اگر احتمال موفقیت تئوریک $p$ باشد، تعداد موفقیت‌های مشاهده‌شده از توزیع دودویی \lr{Binomial(N, p)} پیروی می‌کند. در حد $N \to \infty$، این توزیع به توزیع نرمال $\mathcal{N}(Np, Np(1-p))$ میل می‌کند (\lr{De Moivre-Laplace theorem}). در نتیجه، متغیر استانداردشده $Z$ از توزیع نرمال استاندارد پیروی می‌کند:

\begin{equation}
	Z = \frac{|k - Np|}{\sqrt{Np(1-p)}}
\end{equation}

در این رابطه، $k - Np$ انحراف مشاهده از امید ریاضی است و $\sqrt{Np(1-p)}$ انحراف معیار توزیع دودویی است. اگر $Z$ از یک آستانه بحرانی $Z_{\text{crit}}$ تجاوز کند، فرض صفر (اینکه مشاهده با تئوری هم‌خوان است) رد می‌شود. در سطح معنی‌داری $\alpha = 10^{-4}$، آستانه بحرانی (دوسویه) برابر با $Z_{\text{crit}} \approx 3.8906$ است. این مقدار از تابع توزیع نرمال استاندارد به دست می‌آید:

\begin{equation}
	\alpha = 2 \cdot \Phi(-Z_{\text{crit}}), \quad Z_{\text{crit}} = -\Phi^{-1}(\alpha/2)
\end{equation}

که در آن $\Phi$ تابع توزیع تجمعی نرمال استاندارد است. انتخاب $\alpha = 10^{-4}$ به‌جای $\alpha = 0.05$ (که متداول در علوم اجتماعی است) دو دلیل دارد. نخست، در شبیه‌سازی‌های فیزیکی، ما کنترل کاملی روی منبع تصادفی داریم و می‌توانیم $N$ را به اندازه کافی بزرگ انتخاب کنیم تا نویز آماری کوچک باشد؛ بنابراین آستانه سخت‌گیرانه قابل قبول است. دوم، در \lr{QKD} حتی انحرافات کوچک از تئوری می‌توانند نشانگر باگ‌های ظریف در پیاده‌سازی باشند که در نهایت به نتایج ناامن منجر می‌شوند؛ بنابراین حساسیت بالا ضروری است.

\section{تحلیل معماری و ساختار داده‌ها}

\subsection{ساختار کلی ماژول}

ساختار این ماژول به‌صورت یک \lr{stateless utility module} طراحی شده است؛ یعنی هیچ کلاسی تعریف نشده و تمام توابع به‌صورت مستقل (\lr{pure functions} یا \lr{near-pure functions}) عمل می‌کنند. این طراحی به این دلیل انتخاب شده که اعتبارسنجی فیزیکی باید کاملاً قطعی و قابل تکرار باشد و نباید به وضعیت داخلی هیچ شیء وابسته باشد. هر تابع با دریافت ورودی‌های صریح، خروجی مشخصی تولید می‌کند و عوارض جانبی (\lr{side effects}) آن به چاپ نتیجه آزمون محدود است.

واردسازی (\lr{imports}) در ابتدای ماژول شامل چهار کتابخانه استاندارد پایتون است: \lr{\texttt{math}} برای توابع ریاضی پایه (نظیر \lr{sqrt}، \lr{exp})، \lr{\texttt{sys}} برای دسترسی به \lr{argv} (اگرچه در نسخه فعلی استفاده نمی‌شود، برای توسعه‌های آینده نگه داشته شده)، \lr{\texttt{copy}} برای کپی عمیق پیکربندی، و \lr{\texttt{csv}} برای خواندن فایل خروجی شبیه‌سازی. این حداقل وابستگی تضمین می‌کند که ماژول در هر محیطی قابل اجرا باشد و نیازی به کتابخانه‌های خارجی ندارد.

\begin{LTR}
	\begin{lstlisting}[language=Python]
		import math
		import sys
		import copy
		import csv
	\end{lstlisting}
\end{LTR}

تنها وابستگی خارجی ماژول، به \lr{\texttt{main\_optimized.py}} است که به‌صورت محلی (\lr{local import}) در داخل تابع \lr{run\_isolated\_simulation} بارگذاری می‌شود. این الگو (\lr{lazy import}) دو مزیت دارد: نخست، وابستگی چرخه‌ای (\lr{circular dependency}) را که ممکن است در زمان بارگذاری ماژول رخ دهد، اجتناب می‌کند؛ دوم، اجازه می‌دهد ماژول حتی در صورتی که \lr{\texttt{main\_optimized.py}} در دسترس نباشد، حداقل برای توابع تحلیلی (\lr{channel\_transmittance}، \lr{Q\_mu\_poisson} و غیره) قابل استفاده باشد.

\subsection{توابع تحلیلی channel\_transmittance}

این تابع ساده‌ترین واحد محاسباتی ماژول است و ضریب انتقال کانال را بر اساس قانون \lr{Beer-Lambert} محاسبه می‌کند. با وجود سادگی، این تابع به‌عنوان \lr{building block} پایه برای تمام محاسبات بعدی عمل می‌کند و هر خطا در آن به خطاهای فاحش در تمام آزمون‌ها منجر می‌شود.

\begin{LTR}
	\begin{lstlisting}[language=Python]
		def channel_transmittance(L, alpha_dB_km):
		return 10 ** (-alpha_dB_km * L / 10)
	\end{lstlisting}
\end{LTR}

پارامترهای این تابع عبارت‌اند از: $L$ فاصله به کیلومتر و $\alpha_{\text{dB/km}}$ ضریب تلفت فیبر به \lr{dB/km}. خروجی، ضریب انتقال خطی (در محدوده $[0, 1]$) است. این تابع به‌صورت \lr{stateless} عمل می‌کند و هیچ متغیر سراسری را تغییر نمی‌دهد؛ این ویژگی آن را به یک \lr{pure function} تبدیل می‌کند که آزمون واحد (\lr{unit test}) برای آن بسیار ساده است.

\subsection{توابع تحلیلی Q\_mu\_poisson و Q\_mu\_thermal}

این دو تابع، \lr{Gain} (نرخ کلیک) را برای دو نوع منبع نوری متفاوت محاسبه می‌کنند. تابع \lr{Q\_mu\_poisson} برای منبع \lr{WCP} با توزیع \lr{Poisson} استفاده می‌شود که رایج‌ترین منبع در سیستم‌های \lr{QKD} عملی است.

\begin{LTR}
	\begin{lstlisting}[language=Python]
		def Q_mu_poisson(mu, eta, Y0):
		return 1 - (1 - Y0) * math.exp(-mu * eta)
		
		def Q_mu_thermal(mu, eta, Y0, M=1):
		return 1 - (1 - Y0) / ((1 + mu * eta / M) ** M)
	\end{lstlisting}
\end{LTR}

پارامترهای مشترک این دو تابع عبارت‌اند از: $\mu$ میانگین تعداد فوتون در هر پالس، $\eta$ ضریب انتقال کلی (شامل بازده آشکارساز) و $Y_0$ نرخ دارک کانت. تابع \lr{thermal} پارامتر اضافی $M$ را نیز می‌پذیرد که تعداد حالت‌های مستقل منبع حرارتی را نشان می‌دهد و مقدار پیش‌فرض آن $1$ (تک‌حالت) است.

طراحی این دو تابع به‌گونه‌ای است که هر دو \lr{signature} مشابه دارند (با تفاوت تنها در پارامتر $M$)؛ این امکان را فراهم می‌سازد که در کد فراخوانی، به‌سادگی بین دو مدل منبع جابه‌جا شوند. این الگو (\lr{Strategy pattern}) در مهندسی نرم‌افزار برای پیاده‌سازی الگوریتم‌های قابل تعویض به کار می‌رود.

نکته مهم فیزیکی: در حد $M \to \infty$، تابع \lr{thermal} به \lr{Poisson} میل می‌کند. این را می‌توان با استفاده از حد معروف $\lim_{M \to \infty} (1 + x/M)^M = e^x$ نشان داد:

\begin{equation}
	\lim_{M \to \infty} \frac{1 - Y_0}{\left(1 + \frac{\mu \eta}{M}\right)^M} = (1 - Y_0) \cdot e^{-\mu \eta}
\end{equation}

بنابراین، در حد $M \to \infty$، $Q_\mu^{\text{Thermal}} \to Q_\mu^{\text{Poisson}}$. این ویژگی ریاضی تضمین می‌کند که دو تابع در حد فیزیکی مناسب همگرا می‌شوند و پیاده‌سازی سازگار است.

\subsection{تابع تحلیلی E\_mu\_poisson}

این تابع \lr{QBER} متوسط را برای منبع \lr{WCP} محاسبه می‌کند و بر اساس مدل کلاسیک \lr{GLLP} پیاده‌سازی شده است.

\begin{LTR}
	\begin{lstlisting}[language=Python]
		def E_mu_poisson(mu, eta, Y0, e_d):
		Q = Q_mu_poisson(mu, eta, Y0)
		if Q <= 0: return 0.5
		return (0.5 * Y0 + e_d * (Q - Y0)) / Q
	\end{lstlisting}
\end{LTR}

این تابع ابتدا \lr{Gain} $Q$ را با فراخوانی \lr{Q\_mu\_poisson} محاسبه می‌کند، سپس فرمول \lr{QBER} را اعمال می‌نماید. خط \lr{if Q <= 0: return 0.5} یک محافظ عددی است که در صورتی که $Q$ صفر یا منفی شود (که فیزیکی نیست اما ممکن است در موارد لبه‌ای عددی رخ دهد)، مقدار $0.5$ را برمی‌گرداند. این مقدار $0.5$ نشان‌دهنده آن است که در غیاب هرگونه کلیک، هیچ اطلاعاتی در دست نیست و در نتیجه \lr{QBER} به‌صورت \emph{maximum uncertainty} یعنی $0.5$ تنظیم می‌شود.

پارامتر $e_d$ در این تابع، خطای نوری ذاتی (\lr{misalignment error}) است که از نقص‌های تداخل‌سنج ناشی می‌شود. در عمل، این مقدار معمولاً در بازه $0.005$ تا $0.03$ است و در این ماژول مقدار پیش‌فرض $e_d = 0.01$ به کار رفته است.

\subsection{تابع تحلیلی dead\_time\_correction}

این تابع تصحیح \lr{dead time} را بر اساس مدل \lr{Non-Paralyzable} اعمال می‌کند. این تابع به‌عنوان \lr{building block} پایه برای محاسبه \lr{Gain} در حضور \lr{dead time} عمل می‌کند.

\begin{LTR}
	\begin{lstlisting}[language=Python]
		def dead_time_correction(Q_ideal, tau_d_ns, T_pulse_ns):
		return Q_ideal / (1 + Q_ideal * (tau_d_ns / T_pulse_ns))
	\end{lstlisting}
\end{LTR}

پارامترهای این تابع عبارت‌اند از: $Q_{\text{ideal}}$ احتمال کلیک ایده‌آل (بدون \lr{dead time})، $\tau_d$ زمان مرده به نانوثانیه و $T_{\text{pulse}}$ بازه زمانی میان پالس‌ها به نانوثانیه. خروجی، احتمال کلیک مشاهده‌شده پس از اعمال اثر \lr{dead time} است.

نکته طراحی مهم: این تابع \emph{به‌صورت مستقل} روی یک آشکارساز عمل می‌کند، نه روی کل سیستم. این به این معناست که اگر چندین آشکارساز مستقل وجود داشته باشد (مانند \lr{D0} و \lr{D1} در \lr{BB84})، هر کدام باید به‌صورت جداگانه تصحیح شوند. این موضوع در بخش \lr{TEST 3} به‌صورت صریح پیاده‌سازی شده است.

\subsection{تابع آزمون آماری validate\_statistic}

این تابع هسته آماری ماژول است و آزمون \lr{Z-test} دودویی را روی یک مشاهده انجام می‌دهد. این تابع با دریافت تعداد مشاهده‌شده موفقیت‌ها، تعداد کل آزمایش‌ها، احتمال تئوریک و نام آزمون، نتیجه \lr{PASS}/\lr{FAIL} را برمی‌گرداند.

\begin{LTR}
	\begin{lstlisting}[language=Python]
		def validate_statistic(observed_count, total_trials, expected_prob, test_name):
		expected_count = total_trials * expected_prob
		variance = total_trials * expected_prob * (1 - expected_prob)
		if variance <= 0: variance = 1e-10
		std_dev = math.sqrt(variance)
		z_score = abs(observed_count - expected_count) / std_dev
		
		passed = z_score <= 3.8906
		print(f"{test_name:25} | Exp_p: {expected_prob:.6e} | "
		f"Sim_p: {observed_count/total_trials:.6e} | "
		f"Z: {z_score:6.2f} | {'PASS' if passed else 'FAIL'}")
		return passed
	\end{lstlisting}
\end{LTR}

تحلیل خط‌به‌خط این تابع به شرح زیر است. در خط اول، تعداد مورد انتظار موفقیت‌ها از ضرب تعداد کل آزمایش‌ها در احتمال تئوریک محاسبه می‌شود: $k_{\text{exp}} = N \cdot p$. در خط دوم، واریانس توزیع دودویی از فرمول $\text{Var} = N p (1-p)$ محاسبه می‌شود. در خط سوم، یک محافظ عددی برای جلوگیری از تقسیم بر صفر اعمال می‌شود: اگر واریانس صفر یا منفی باشد (که ممکن است در موارد لبه‌ای مانند $p = 0$ یا $p = 1$ رخ دهد)، مقدار بسیار کوچک $10^{-10}$ جایگزین می‌شود. در خط چهارم، انحراف معیار از جذر واریانس محاسبه می‌شود. در خط پنجم، \lr{Z-score} از نسبت قدرمطلق انحراف به انحراف معیار به دست می‌آید. در خط هفتم، نتیجه آزمون با مقایسه \lr{Z-score} با آستانه بحرانی $Z_{\text{crit}} = 3.8906$ (متناظر با $\alpha = 10^{-4}$ دوسویه) تعیین می‌شود.

قابل توجه است که این تابع در خط چاپ، علاوه بر نتیجه \lr{PASS}/\lr{FAIL}، مقادیر احتمال تئوریک، احتمال شبیه‌سازی‌شده و \lr{Z-score} را نیز نمایش می‌دهد. این اطلاعات تشخیصی (\lr{diagnostic information}) برای تحلیل شکست آزمون‌ها بسیار ارزشمند است؛ زیرا به توسعه‌دهنده اجازه می‌دهد ببیند انحراف از تئوری چقدر بزرگ است و در چه جهتی (مثبت یا منفی) قرار دارد. این الگو (\lr{verbose assertion}) در توسعه نرم‌افزار علمی به‌عنوان بهترین روش شناخته می‌شود.

\subsection{تابع اجرای شبیه‌سازی run\_isolated\_simulation}

این تابع پیچیده‌ترین بخش ماژول است و یک نشست شبیه‌سازی مستقل روی موتور \lr{\texttt{main\_optimized.py}} اجرا می‌کند. این تابع با دریافت یک دیکشنری \lr{config\_overrides}، پیکربندی پیش‌فرض را با مقادیر دلخواه بازنویسی می‌کند و شبیه‌سازی را با پارامترهای ایزوله اجرا می‌نماید.

\begin{LTR}
	\begin{lstlisting}[language=Python]
		def run_isolated_simulation(config_overrides):
		import main_optimized as m
		cfg = copy.deepcopy(m.DEFAULT_CONFIG)
		
		N_pulses = 100000
		dist = config_overrides.get("dist", 0)
		
		cfg["simulation"]["apply_statistical_noise"] = True
		cfg["simulation"]["distance_start_km"] = dist
		cfg["simulation"]["distance_stop_km"] = dist
		cfg["simulation"]["distance_points"] = 1
		cfg["simulation"]["min_pulses_log"] = 5
		cfg["simulation"]["max_pulses_log"] = 5
		
		mu = config_overrides.get("mu", 0.5)
		cfg["source"]["statistics_type"] = config_overrides.get("stat_type", "POISSON")
		cfg["source"]["intensity_jitter"] = 0.0
		
		cfg["source"]["pulses"] = {
			"signal": {"mu": mu, "prob": 1.0},
			"decoy": {"mu": 0.1, "prob": 0.0},
			"vacuum": {"mu": 0.0, "prob": 0.0}
		}
		
		eta_det = config_overrides.get("det_eff", 0.15)
		Y0 = config_overrides.get("Y0", 3e-6)
		ed = config_overrides.get("ed", 0.01)
		
		cfg["detector"]["det_eff_d0"] = eta_det
		cfg["detector"]["det_eff_d1"] = eta_det
		cfg["detector"]["dark_rate"] = Y0
		cfg["detector"]["dark_rate_d1"] = Y0
		cfg["detector"]["qber_intrinsic"] = ed
		cfg["detector"]["misalignment"] = 0.0
		
		cfg["detector"]["dead_time_ns"] = config_overrides.get("dt", 0.0)
		cfg["detector"]["dead_time_model"] = "NON_PARALYZABLE"
		
		m.run_and_save_csv(cfg, num_workers=1, sweeps={})
		
		with open("qkd_results_optimized_sweep.csv", "r") as f:
		rows = list(csv.DictReader(f))
		latest = rows[-1]
		
		y_sig = float(latest.get("detection_yield_signal", 0))
		q_sig = float(latest.get("qber_signal", 0))
		t_pulse = float(cfg["source"]["pulse_period_ns"])
		
		return {
			"clicks": int(y_sig * N_pulses),
			"errors": int(q_sig * y_sig * N_pulses),
			"pulse_period": t_pulse
		}
	\end{lstlisting}
\end{LTR}

تحلیل این تابع به شرح زیر است. در خط اول، ماژول \lr{\texttt{main\_optimized}} به‌صورت محلی (\lr{lazy import}) بارگذاری می‌شود؛ این الگو امکان اجرای ماژول \lr{validate\_physics} را حتی در صورتی که \lr{\texttt{main\_optimized}} در دسترس نباشد، برای توابع تحلیلی فراهم می‌سازد. در خط دوم، یک کپی عمیق از پیکربندی پیش‌فرض \lr{DEFAULT\_CONFIG} ساخته می‌شود. استفاده از \lr{deepcopy} ضروری است زیرا پیکربندی یک دیکشنری تودرتو است و کپی سطحی باعث به اشتراک‌گذاری زیردیکشنری‌ها می‌شود که می‌تواند به تغییرات ناخواسته در پیکربندی پیش‌فرض منجر گردد.

در خطوط بعدی، پیکربندی شبیه‌سازی به‌صورت ایزوله تنظیم می‌شود: \lr{apply\_statistical\_noise = True} فعال‌سازی نویز آماری، \lr{distance\_start\_km = distance\_stop\_km = dist} اجرای شبیه‌سازی در یک فاصله منفرد (بدون \lr{sweep})، \lr{distance\_points = 1} تولید تنها یک نقطه داده، و \lr{min\_pulses\_log = max\_pulses\_log = 5} تنظیم تعداد پالس‌ها به $10^5 = 100000$.

یکی از نکات کلیدی این تابع، تنظیم \lr{prob = 1.0} برای سیگنال و \lr{prob = 0.0} برای \lr{decoy} و \lr{vacuum} است. این تنظیم به این معناست که \emph{تمام} پالس‌ها به‌عنوان سیگنال ارسال می‌شوند و هیچ پالس \lr{decoy} یا \lr{vacuum} وجود ندارد. این ایزوله‌سازی ضروری است زیرا معادلات تحلیلی مرجع (مانند \lr{Q\_mu\_poisson}) برای پالس سیگنال خالص به‌دست آمده‌اند و حضور \lr{decoy} یا \lr{vacuum} نیاز به تجمیع متفاوت دارد.

در بخش آشکارساز، چهار پارامتر کلیدی تنظیم می‌شوند: \lr{det\_eff\_d0} و \lr{det\_eff\_d1} بازده آشکارساز \lr{D0} و \lr{D1}، \lr{dark\_rate} و \lr{dark\_rate\_d1} نرخ دارک کانت هر آشکارساز، \lr{qber\_intrinsic} خطای ذاتی، و \lr{misalignment = 0.0} حذف خطای \lr{misalignment} (زیرا این کمیت به‌صورت جداگانه در پارامتر $e_d$ کنترل می‌شود).

در خط فراخوانی \lr{run\_and\_save\_csv}، شبیه‌سازی با پیکربندی نهایی اجرا می‌شود و خروجی در فایل \lr{qkd\_results\_optimized\_swap.csv} ذخیره می‌گردد. سپس این فایل خوانده شده و آخرین ردیف (که نتیجه شبیه‌سازی است) استخراج می‌شود. در نهایت، سه مقدار کلیدی برگردانده می‌شود: تعداد کلیک‌ها (از ضرب \lr{detection\_yield\_signal} در تعداد پالس‌ها)، تعداد خطاها (از ضرب \lr{QBER} در کلیک‌ها)، و بازه زمانی پالس (برای محاسبه \lr{dead time}).

\section{پیاده‌سازی ریاضی و الگوریتمی}

\subsection{الگوریتم TEST 1: اعتبارسنجی منبع Poisson}

این آزمون، منبع \lr{WCP} با توزیع \lr{Poisson} را در سه فاصله $0$، $50$ و $100$ کیلومتر اعتبارسنجی می‌کند. برای هر فاصله، دو کمیت مستقل آزمون می‌شود: \lr{Gain} (نرخ کلیک) و \lr{QBER} (نرخ خطای بیت).

\begin{LTR}
	\begin{lstlisting}[language=Python]
		print("\n[TEST 1] Poisson WCP Asymptotic")
		for d in [0, 50, 100]:
		res = run_isolated_simulation({
			"dist": d, "mu": 0.5, "stat_type": "POISSON",
			"det_eff": eta_det, "Y0": Y0, "ed": ed, "dt": 0.0
		})
		eta = channel_transmittance(d, 0.2) * eta_det
		Q_th = Q_mu_poisson(0.5, eta, Y0) * 0.5
		E_th = E_mu_poisson(0.5, eta, Y0, ed)
		validate_statistic(res["clicks"], N, Q_th, f"WCP Gain L={d}km")
		validate_statistic(res["errors"], res["clicks"], E_th, f"WCP QBER L={d}km")
	\end{lstlisting}
\end{LTR}

تحلیل ریاضی این آزمون به شرح زیر است. در هر تکرار، ابتدا شبیه‌سازی با پارامترهای ایزوله اجرا می‌شود و تعداد کلیک‌ها و خطاها به دست می‌آید. سپس ضریب انتقال تئوریک از ضرب ضریب انتقال کانال (با $\alpha_{\text{dB/km}} = 0.2$) در بازده آشکارساز محاسبه می‌شود:

\begin{equation}
	\eta_{\text{total}}(d) = \eta_{\text{det}} \cdot 10^{-\frac{0.2 \cdot d}{10}}
\end{equation}

برای $d = 0$ کیلومتر: $\eta_{\text{total}} = 0.15$، برای $d = 50$ کیلومتر: $\eta_{\text{total}} = 0.015$، و برای $d = 100$ کیلومتر: $\eta_{\text{total}} = 0.0015$.

سپس \lr{Gain} تئوریک از فرمول \lr{Q\_mu\_poisson} محاسبه می‌شود و در $0.5$ ضرب می‌گردد تا اثر \lr{sifting} اعمال شود. این ضرب در $0.5$ به این دلیل است که در \lr{BB84}، گیرنده به‌صورت تصادفی یکی از دو پایه را انتخاب می‌کند و تنها در نیمی از موارد، پایه گیرنده با فرستنده هم‌خوان است (\lr{sifting}). در نتیجه، نرخ کلیک مشاهده‌شده نیمی از \lr{Gain} خام است.

\lr{QBER} تئوریک از فرمول \lr{E\_mu\_poisson} به دست می‌آید. نکته مهم آن است که در این محاسبه، \lr{sifting} اعمال نمی‌شود زیرا \lr{QBER} نسبت خطا به کل کلیک‌ها است و این نسبت تحت \lr{sifting} تغییر نمی‌کند (هم خطاها و هم کل کلیک‌ها به همان نسبت $0.5$ کاهش می‌یابند).

دو آزمون آماری جداگانه انجام می‌شود: یکی برای \lr{Gain} (با $N = 100000$ آزمایش و احتمال تئوریک $Q_{\text{th}}$) و دیگری برای \lr{QBER} (با $N = $ تعداد کلیک‌ها و احتمال تئوریک $E_{\text{th}}$). در آزمون \lr{QBER}، تعداد آزمایش‌ها به تعداد کلیک‌ها تغییر می‌کند زیرا \lr{QBER} بر روی کلیک‌ها اعمال می‌شود، نه بر روی کل پالس‌ها. این تمایز ظریف برای صحت آماری آزمون ضروری است.

\subsection{الگوریتم TEST 2: اعتبارسنجی منبع Thermal}

این آزمون، منبع حرارتی تک‌حالت را در فاصله صفر اعتبارسنجی می‌کند. این آزمون کوتاه‌تر از \lr{TEST 1} است زیرا فقط \lr{Gain} بررسی می‌شود و فقط در یک فاصله.

\begin{LTR}
	\begin{lstlisting}[language=Python]
		print("\n[TEST 2] Thermal Single-Mode")
		res = run_isolated_simulation({
			"dist": 0, "mu": 0.5, "stat_type": "THERMAL",
			"det_eff": eta_det, "Y0": Y0, "dt": 0.0
		})
		Q_th = Q_mu_thermal(0.5, eta_det, Y0, M=1) * 0.5
		validate_statistic(res["clicks"], N, Q_th, "Thermal Gain L=0km")
	\end{lstlisting}
\end{LTR}

در این آزمون، پارامتر \lr{stat\_type = "THERMAL"} به موتور شبیه‌سازی دستور می‌دهد که به‌جای توزیع \lr{Poisson}، از توزیع \lr{Bose-Einstein} تک‌حالت برای تولید فوتون استفاده کند. سپس \lr{Gain} تئوریک از فرمول \lr{Q\_mu\_thermal} با $M = 1$ محاسبه می‌شود:

\begin{equation}
	Q_\mu^{\text{Thermal}}(\mu = 0.5, \eta = 0.15, Y_0 = 3 \times 10^{-6}, M = 1) = 1 - \frac{1 - 3 \times 10^{-6}}{1 + 0.5 \times 0.15} \approx 0.075
\end{equation}

این مقدار در $0.5$ ضرب می‌شود تا \lr{sifting} اعمال شود و در نتیجه $Q_{\text{th}} \approx 0.0375$ به دست می‌آید. این یعنی از $100000$ پالس، انتظار می‌رود حدود $3750$ کلیک ثبت شود.

نکته فیزیکی قابل توجه: برای همان $\mu$ و $\eta$، \lr{Gain} منبع حرارتی بزرگ‌تر از \lr{Gain} منبع \lr{Poisson} است. این به دلیل پدیده \lr{bunching} در نور حرارتی است که احتمال پالس‌های چندفوتونی را افزایش می‌دهد. در مقابل، نور \lr{Poisson} دارای آمار \lr{coherent} است که در آن فوتون‌ها مستقل توزیع می‌شوند. این تفاوت آماری یکی از مفاهیم بنیادین اپتیک کوانتومی است که در آزمون ما به‌صورت کمی نشان داده می‌شود.

\subsection{الگوریتم TEST 3: اعتبارسنجی Dead Time مستقل SPAD}

این آزمون پیچیده‌ترین آزمون ماژول است و فیزیک \lr{dead time} در حضور دو آشکارساز مستقل را بررسی می‌کند. این آزمون به‌طور خاص این موضوع مهم را راستی‌آزمایی می‌کند که هر \lr{SPAD} به‌طور مستقل زمان مرده خود را تجربه می‌کند، نه اینکه زمان مرده روی کل سیستم اعمال شود.

\begin{LTR}
	\begin{lstlisting}[language=Python]
		print("\n[TEST 3] Dead-Time Physics (Independent SPADs)")
		dt_ns = 50.0
		res = run_isolated_simulation({
			"dist": 0, "mu": 0.5, "stat_type": "POISSON",
			"det_eff": eta_det, "Y0": Y0, "dt": dt_ns
		})
		
		# فیزیک صحیح: محاسبه زمان مرده روی هر SPAD به طور مستقل
		Q_ideal_raw = Q_mu_poisson(0.5, eta_det, Y0)
		Q_per_spad = Q_ideal_raw / 2.0
		T_pulse = res["pulse_period"]
		
		Q_dead_per_spad = dead_time_correction(Q_per_spad, dt_ns, T_pulse)
		Q_dt_exp = Q_dead_per_spad * 2.0 * 0.5  # 2 SPADs * Sifting probability
		
		validate_statistic(res["clicks"], N, Q_dt_exp, f"Dead-time tau={dt_ns}ns")
	\end{lstlisting}
\end{LTR}

تحلیل ریاضی این آزمون به شرح زیر است. ابتدا \lr{Gain} خام ایده‌آل (بدون \lr{dead time}) از فرمول \lr{Q\_mu\_poisson} محاسبه می‌شود. سپس این مقدار به دو نیمه تقسیم می‌شود زیرا در سیستم \lr{BB84}، هر کلیک به‌صورت تصادفی به یکی از دو آشکارساز \lr{D0} یا \lr{D1} می‌رود (با احتمال برابر $0.5$). در نتیجه، \lr{Gain} هر آشکارساز برابر با $Q_{\text{per-spad}} = Q_{\text{ideal}}^{\text{raw}} / 2$ است.

سپس تصحیح \lr{dead time} به‌صورت مستقل روی هر آشکارساز اعمال می‌شود:

\begin{equation}
	Q_{\text{obs}}^{\text{per-spad}} = \frac{Q_{\text{per-spad}}}{1 + Q_{\text{per-spad}} \cdot \frac{\tau_d}{T_{\text{pulse}}}}
\end{equation}

در نهایت، برای محاسبه \lr{Gain} کل سیستم، مقدار هر آشکارساز در $2$ ضرب می‌شود (چون دو آشکارساز مستقل وجود دارد) و سپس در $0.5$ ضرب می‌شود تا اثر \lr{sifting} اعمال گردد:

\begin{equation}
	Q_{\text{final}} = Q_{\text{obs}}^{\text{per-spad}} \cdot 2 \cdot 0.5 = Q_{\text{obs}}^{\text{per-spad}}
\end{equation}

این رابطه جالب نشان می‌دهد که در نهایت، \lr{Gain} مشاهده‌شده با \lr{Gain} هر آشکارساز (پس از تصحیح \lr{dead time}) برابر است. این ساده‌سازی ریاضی به این دلیل رخ می‌دهد که عوامل $2$ (برای دو آشکارساز) و $0.5$ (برای \lr{sifting}) یکدیگر را خنثی می‌کنند. با این حال، \emph{ترتیب اعمال} این عوامل مهم است: ابتدا باید \lr{dead time} روی هر آشکارساز به‌صورت مستقل اعمال شود و سپس جمع انجام گردد. اگر به‌اشتباه \lr{dead time} روی \lr{Gain} کل اعمال شود، نتیجه متفاوت و نادرستی به دست می‌آید.

برای درک کمی، فرض کنید $Q_{\text{ideal}}^{\text{raw}} = 0.075$ (مقدار تقریبی برای $\mu = 0.5$، $\eta = 0.15$، $Y_0 = 3 \times 10^{-6}$ در فاصله صفر). آنگاه $Q_{\text{per-spad}} = 0.0375$. اگر $\tau_d = 50$ نانوثانیه و $T_{\text{pulse}} = 1000$ نانوثانیه (نرخ پالس $1$ مگاهرتز) باشد:

\begin{equation}
	Q_{\text{obs}}^{\text{per-spad}} = \frac{0.0375}{1 + 0.0375 \times \frac{50}{1000}} = \frac{0.0375}{1.001875} \approx 0.03743
\end{equation}

این نشان می‌دهد که در این شرایط، \lr{dead time} باعث کاهش $\approx 0.18\%$ در \lr{Gain} می‌شود. هرچند این کاهش کوچک به نظر می‌رسد، اما در آزمون آماری با $N = 100000$ قابل شناسایی است زیرا انحراف معیار تئوریک برابر با $\sqrt{100000 \times 0.0375 \times 0.9625} \approx 60$ است و انحراف $\approx 0.18\%$ از $0.0375$ معادل $\approx 6.75$ کلیک است که بسیار کوچک‌تر از $3.89 \times 60 \approx 233$ (\lr{Z-crit}) است. اما اگر \lr{dead time} به‌اشتباه روی کل \lr{Gain} اعمال شود، خطا به‌مراتب بزرگ‌تر خواهد بود و آزمون شکست می‌خورد.

\subsection{محاسبه Q-score و آستانه بحرانی}

انتخاب آستانه بحرانی $Z_{\text{crit}} = 3.8906$ بر اساس سطح معنی‌داری $\alpha = 10^{-4}$ دوسویه است. این مقدار از معادله زیر به دست می‌آید:

\begin{equation}
	\alpha = P(|Z| > Z_{\text{crit}} | H_0) = 2 \cdot \Phi(-Z_{\text{crit}})
\end{equation}

که در آن $\Phi$ تابع توزیع تجمعی نرمال استاندارد است. با حل معادله $\Phi(-Z_{\text{crit}}) = 5 \times 10^{-5}$، مقدار $Z_{\text{crit}} \approx 3.8906$ به دست می‌آید. این مقدار در کد به‌صورت یک ثابت سخت‌کد شده (\lr{hard-coded}) تعریف شده است.

در عمل، انتخاب این سطح سخت‌گیرانه دو دلیل دارد. نخست، در شبیه‌سازی \lr{QKD} با $N = 100000$ پالس، انحراف معیار نسبی برابر با $\sqrt{p(1-p)/N} \approx 0.0016$ برای $p \approx 0.04$ است. این یعنی نویز آماری ذاتی شبیه‌سازی در حد $0.16\%$ است و انحرافات بزرگ‌تر از این مقدار باید شناسایی شوند. با $Z_{\text{crit}} = 3.89$، آستانه شناسایی برابر با $3.89 \times 0.16\% \approx 0.62\%$ است که به‌طور معناداری بزرگ‌تر از نویز است. دوم، این سطح به‌قدری سخت‌گیرانه است که \lr{false positive} بسیار نادر است (احتمال $\alpha = 10^{-4}$ یعنی یک شکست در هر $10000$ اجرای مستقل) اما هنوز به اندازه کافی حساس است که باگ‌های واقعی را شناسایی کند.

\subsection{تحلیل پیچیدگی و کارایی}

از نظر پیچیدگی زمانی، تمام توابع تحلیلی (\lr{channel\_transmittance}، \lr{Q\_mu\_poisson}، \lr{Q\_mu\_thermal}، \lr{E\_mu\_poisson}، \lr{dead\_time\_correction}) در زمان $\mathcal{O}(1)$ اجرا می‌شوند زیرا تنها شامل چند عملیات حسابی ساده هستند. تابع \lr{validate\_statistic} نیز در زمان $\mathcal{O}(1)$ اجرا می‌شود.

تابع \lr{run\_isolated\_simulation} پیچیده‌تر است: زمان اجرای آن به زمان اجرای \lr{run\_and\_save\_csv} در \lr{\texttt{main\_optimized.py}} بستگی دارد که خود به تعداد پالس‌ها $N$ و پیچیدگی خط لوله شبیه‌سازی وابسته است. برای $N = 100000$ پالس، این زمان معمولاً در حد چند ثانیه تا چند ده ثانیه است. خواندن فایل \lr{CSV} در زمان $\mathcal{O}(R)$ اجرا می‌شود که در آن $R$ تعداد ردیف‌های فایل است (که در این حالت تنها $1$ ردیف است زیرا \lr{sweep} خالی است).

از نظر پیچیدگی حافظه، ماژول بسیار سبک است؛ تنها ساختار داده‌های قابل توجه دیکشنری \lr{cfg} (با اندازه ثابت) و لیست \lr{rows} (با اندازه $1$) هستند. در نتیجه، کل ماژول می‌تواند به‌راحتی در محیط‌های با حافظه محدود اجرا شود.

\section{ارسال و دریافت با سایر ماژول‌ها}

\subsection{وابستگی به main\_optimized.py}

مهم‌ترین وابستگی خارجی این ماژول، به فایل \lr{\texttt{main\_optimized.py}} است که به‌عنوان موتور اصلی شبیه‌سازی عمل می‌کند. این وابستگی از طریق سه مؤلفه کلیدی برقرار می‌شود: \lr{DEFAULT\_CONFIG} به‌عنوان پیکربندی پایه، \lr{run\_and\_save\_csv} به‌عنوان نقطه ورود شبیه‌سازی، و فرمت فایل \lr{CSV} خروجی به‌عنوان پروتکل ارتباطی.

\begin{LTR}
	\begin{lstlisting}[language=Python]
		import main_optimized as m
		cfg = copy.deepcopy(m.DEFAULT_CONFIG)
		# ... (modify cfg) ...
		m.run_and_save_csv(cfg, num_workers=1, sweeps={})
	\end{lstlisting}
\end{LTR}

این الگوی وابستگی به این معناست که هرگونه تغییر در \lr{DEFAULT\_CONFIG} (مانند افزودن فیلدهای جدید یا تغییر نام فیلدها) یا در فرمت خروجی \lr{CSV} (مانند تغییر نام ستون‌ها) می‌تواند به شکست ماژول \lr{validate\_physics} منجر شود. برای کاهش این ریسک، ماژول از کلیدهای رشته‌ای (\lr{string keys}) با مقادیر پیش‌فرض استفاده می‌کند (مانند \lr{\texttt{latest.get("detection\_yield\_signal", 0)}}) که در صورت نبود یک کلید، مقدار پیش‌فرض برمی‌گردد و از خطای \lr{KeyError} جلوگیری می‌کند.

\subsection{ورودی‌های ماژول}

ورودی اصلی ماژول از طریق دیکشنری \lr{config\_overrides} در تابع \lr{run\_isolated\_simulation} ارائه می‌شود. این دیکشنری شامل هفت کلید اختیاری است: \lr{dist} (فاصله به کیلومتر، پیش‌فرض $0$)، \lr{mu} (میانگین فوتون، پیش‌فرض $0.5$)، \lr{stat\_type} (نوع توزیع آماری، پیش‌فرض \lr{"POISSON"})، \lr{det\_eff} (بازده آشکارساز، پیش‌فرض $0.15$)، \lr{Y0} (نرخ دارک کانت، پیش‌فرض $3 \times 10^{-6}$)، \lr{ed} (خطای ذاتی، پیش‌فرض $0.01$) و \lr{dt} (زمان مرده به نانوثانیه، پیش‌فرض $0$).

این الگو (\lr{configuration override pattern}) مزایای متعددی دارد. نخست، امکان تنظیم دقیق پارامترهای شبیه‌سازی برای هر آزمون به‌صورت مستقل فراهم می‌شود. دوم، مقادیر پیش‌فرض معنی‌دار هستند و در صورت عدم ارائه یک کلید، مقدار معقولی استفاده می‌شود. سوم، افزودن پارامترهای جدید در آینده بدون شکستن کدهای فراخوان موجود امکان‌پذیر است.

در سطح بالاتر، ماژول سه گروه پارامتر ثابت در بدنه اصلی (\lr{\_\_main\_\_}) تعریف می‌کند: $N = 100000$ (تعداد پالس‌ها)، $\eta_{\text{det}} = 0.15$ (بازده آشکارساز)، $Y_0 = 3 \times 10^{-6}$ (نرخ دارک کانت) و $e_d = 0.01$ (خطای ذاتی). این مقادیر به‌عنوان \lr{golden parameters} عمل می‌کنند و در تمام آزمون‌ها به‌صورت ثابت استفاده می‌شوند. ثابت نگه‌داشتن این پارامترها ضروری است زیرا تغییر آن‌ها می‌تواند به نتایج متفاوتی منجر شود که مقایسه با مقادیر تئوریک را دشوار سازد.

\subsection{خروجی‌های ماژول}

خروجی اصلی ماژول به دو صورت است: خروجی متنی در \lr{stdout} و مقدار بازگشتی بولی از تابع \lr{validate\_statistic}. خروجی متنی شامل یک جدول قالب‌بندی‌شده است که برای هر آزمون، نام، احتمال تئوریک، احتمال شبیه‌سازی‌شده، \lr{Z-score} و نتیجه \lr{PASS}/\lr{FAIL} را نمایش می‌دهد. این خروجی برای تحلیل بصری توسط توسعه‌دهنده طراحی شده است.

علاوه بر خروجی متنی، تابع \lr{validate\_statistic} مقدار بولی \lr{passed} را برمی‌گرداند که می‌تواند برای اتوماسیون در \lr{CI/CD pipeline} استفاده شود. در یکپارچه‌سازی مستمر، اسکریپت می‌تواند تمام مقادیر بازگشتی را جمع کرده و در صورت وجود حتی یک \lr{FAIL}، کد خطا برگرداند تا \lr{pipeline} متوقف شود. این الگو (\lr{exit code based gating}) در ابزارهایی مانند \lr{GitHub Actions} و \lr{GitLab CI} به‌صورت استاندارد پشتیبانی می‌شود.

تابع \lr{run\_isolated\_simulation} نیز یک دیکشنری با سه کلید برمی‌گرداند: \lr{clicks} (تعداد کلیک‌های ثبت‌شده)، \lr{errors} (تعداد خطاهای ثبت‌شده) و \lr{pulse\_period} (بازه زمانی پالس به نانوثانیه). این مقادیر به‌عنوان ورودی به توابع آزمون آماری استفاده می‌شوند و امکان محاسبه مقادیر تئوریک متناظر را فراهم می‌سازند.

\subsection{ارتباط با فایل CSV خروجی}

یکی از جنبه‌های مهم ارتباط با موتور شبیه‌سازی، فرمت فایل \lr{CSV} خروجی است. پس از اجرای \lr{run\_and\_save\_csv}، نتایج در فایل \lr{qkd\_results\_optimized\_sweep.csv} ذخیره می‌شوند و ماژول \lr{validate\_physics} این فایل را می‌خواند.

\begin{LTR}
	\begin{lstlisting}[language=Python]
		with open("qkd_results_optimized_sweep.csv", "r") as f:
		rows = list(csv.DictReader(f))
		latest = rows[-1]
		
		y_sig = float(latest.get("detection_yield_signal", 0))
		q_sig = float(latest.get("qber_signal", 0))
		t_pulse = float(cfg["source"]["pulse_period_ns"])
	\end{lstlisting}
\end{LTR}

استفاده از \lr{csv.DictReader} به این معناست که فایل \lr{CSV} باید دارای سطر سرستون باشد و نام ستون‌ها باید دقیقاً با کلیدهای مورد انتظار (\lr{detection\_yield\_signal} و \lr{qber\_signal}) مطابقت داشته باشد. این یک \emph{قرارداد ارتباطی} (\lr{interface contract}) بین دو ماژول است که باید در صورت تغییر فرمت \lr{CSV}، به‌روزرسانی شود.

نکته طراحی مهم: استفاده از \lr{rows[-1]} برای گرفتن آخرین ردیف به این دلیل است که در حالت \lr{sweep} با چندین نقطه فاصله، چندین ردیف تولید می‌شود. در این ماژول، چون \lr{sweep} خالی است و فقط یک فاصله مشخص شده، تنها یک ردیف تولید می‌شود و \lr{rows[-1]} همان ردیف اول است. اما استفاده از \lr{rows[-1]} به‌جای \lr{rows[0]} یک الگوی دفاعی است که در صورتی که به‌اشتباه چندین ردیف تولید شود، آخرین ردیف (که معمولاً متناظر با آخرین فاصله است) را برمی‌گرداند.

نکته دیگر آن است که \lr{pulse\_period} از \lr{cfg} خوانده می‌شود، نه از فایل \lr{CSV}. این به این دلیل است که \lr{pulse\_period} یک پارامتر پیکربندی است که در طول شبیه‌سازی ثابت می‌ماند و نیازی به ذخیره در خروجی ندارد. با این حال، این الگو به یکپارچگی داده‌ها آسیب می‌زند زیرا اگر \lr{cfg} پس از اجرای شبیه‌سازی تغییر کند، مقدار خوانده‌شده ممکن است با مقدار واقعی استفاده‌شده در شبیه‌سازی متفاوت باشد. در یک طراحی بهتر، این مقدار نیز در فایل \lr{CSV} ذخیره می‌شد.

\subsection{یکپارچه‌سازی با سیستم CI/CD}

این ماژول به‌عنوان یک \lr{test suite} مستقل عمل می‌کند و می‌تواند به‌راحتی در سیستم‌های \lr{CI/CD} یکپارچه شود. الگوی متداول استفاده به این صورت است:

\begin{LTR}
	\begin{lstlisting}[language=Python]
		# In CI script:
		python validate_physics.py
		if [ $? -ne 0 ]; then
		echo "Physics validation failed!"
		exit 1
		fi
	\end{lstlisting}
\end{LTR}

با این حال، در نسخه فعلی ماژول، کد خروج (\lr{exit code}) به‌صورت صریح تنظیم نمی‌شود و حتی در صورت شکست آزمون‌ها، اسکریپت با کد $0$ (موفقیت) خارج می‌شود. این یک محدودیت است که در نسخه‌های آینده باید اصلاح شود؛ به‌عنوان مثال، می‌توان در انتهای بدنه اصلی، تعداد شکست‌ها را شمارش کرده و در صورت بزرگ‌تر بودن از صفر، با \lr{sys.exit(1)} خارج شد.

\subsection{محدودیت‌ها و توسعه‌های آینده}

این ماژول دارای چند محدودیت است که در توسعه‌های آینده قابل بهبود است. نخست، تعداد پالس‌ها $N = 100000$ ثابت است و قابل تنظیم نیست؛ برای آزمون‌های حساس‌تر، می‌توان $N$ را به $10^6$ یا $10^7$ افزایش داد، هرچند این به زمان اجرای طولانی‌تر نیاز دارد. دوم، تنها سه آزمون تعریف شده است و سناریوهای پیشرفته‌تری مانند \lr{decoy-state} کامل، \lr{afterpulse}، و \lr{polarization drift} پوشش داده نمی‌شوند. سوم، خروجی به‌صورت متنی است و برای تحلیل خودکار نیاز به تجزیه دارد؛ می‌توان خروجی را به فرمت \lr{JSON} یا \lr{JUnit XML} تغییر داد تا با ابزارهای \lr{CI} بهتر یکپارچه شود. چهارم، آزمون‌ها به‌صورت قطعی (\lr{deterministic}) هستند زیرا از بذر تصادفی ثابت استفاده نمی‌کنند؛ در نتیجه، در اجراهای مختلف ممکن است نتایج کمی متفاوت باشند. برای اصلاح این، می‌توان از \lr{SeedSequence} با بذر ثابت استفاده کرد.

علاوه بر این، ماژول می‌تواند با معرفی \lr{parametrized tests} گسترش یابد؛ به‌عنوان مثال، می‌توان به‌جای سه فاصله ثابت در \lr{TEST 1}، یک \lr{sweep} روی بازه $[0, 100]$ کیلومتر با گام $10$ کیلومتر انجام داد و آزمون را برای هر نقطه تکرار کرد. این الگو (\lr{parameterized testing}) در فریم‌ورک‌هایی مانند \lr{pytest} به‌صورت استاندارد پشتیبانی می‌شود و می‌تواند اعتبار فیزیکی ماژول را در طیف وسیع‌تری از پارامترها تضمین نماید.

در نهایت، این ماژول به‌عنوان یک \lr{golden reference} برای فیزیک \lr{QKD} در فریم‌ورک عمل می‌کند و هرگونه تغییر در پیاده‌سازی موتور فیزیکی باید پیش از \lr{merge} با این آزمون‌ها اعتبارسنجی شود. این رویکرد \lr{shift-left testing} تضمین می‌کند که باگ‌های فیزیکی در همان مراحل اولیه توسعه شناسایی شوند و به انتشار نتایج نادرست در مقالات یا گزارش‌های فنی منجر نشوند. وجود چنین ماژول اعتبارسنجی مستقل، نشانه بلوغ مهندسی نرم‌افزار در پروژه‌های علمی است و در پروژه‌های تحقیقاتی جدی به‌عنوان یک ضرورت در نظر گرفته می‌شود.
